\section{标准、参考架构与方法}

\subsection{标准的重要性}

工业4.0只有依靠统一的规范与标准，收益才能真正被兑现。与此同时，安全本身也是工业4.0一切应用场景得以展开的前提与推动力，安全设计这一原则将会贯穿其中。

工业4.0所追求的互联互通远远超出了单纯的信息传递。必须发展出一套各方共同理解的语言：从符号、数据、信息到知识，每一层级都对应着语法、语义与语用这几重加工，只有把符号组合成有意义的数据、再赋予数据以语境化的语义、进而将信息网络化为可用于解决问题的知识，沟通才能真正达成一致。

从历史的角度看，标准化本身也经历了一个从被动走向主动的转变。早期的标准化更多是对既有产品的事后规范，而后逐渐延伸到对工艺流程的规范，如今则进一步发展为在创新产生的同时就主动展开的标准化工作。

在具体效力层面，规范、标准与事实标准之间存在明显差异，这种差异直接影响着企业所拥有的灵活性。处于金字塔顶端、具有完全共识性质的是各层级的正式标准，包括国际层面的ISO，欧洲层面的CEN，以及国家层面的DIN，其作用在于保障某一领域核心任务的可靠性。居中的是有限共识性质的联盟标准与规范，例如源自工业4.0平台的RAMI与I4.0组件，其作用在于持续提升互操作性。处于底部、由市场与企业自身渗透形成的则是事实标准，例如企业标准与厂内规范（例如Windows系统），其特点是灵活性最大。

\subsection{两大参考架构模型：RAMI 4.0与IIRA}

面对工业4.0所涉及的庞杂议题，如何以一种结构化的方式加以把握，是参考架构模型试图回答的问题。德国的RAMI4.0与美国的IIRA是两个彼此独立发展、却又出发点各异的代表性方案：工业4.0平台由德国 BITKOM、VDMA、ZVEI 三大行业协会于2013年共同发起，并于2014年四月发布了首个RAMI4.0。与此同时，OMG 于2014年为工业互联网联盟 IIC 建立平台，其创始成员包括IBM、Cisco、AT\&T、Intel 和 GE，并于2015年推出IIRA。RAMI4.0源自一项由政府主导的工业4.0平台倡议，而IIRA则完全由产业界自发推动。就覆盖领域而言，IIRA的视角横跨能源、医疗卫生、制造、公共部门与交通运输，进而汇聚到工业4.0所对应的制造领域，其定位更偏向纯粹的信息技术架构；而RAMI4.0则聚焦生产制造，对下一代制造价值链给出更为详尽的模型。

两个架构模型彼此并不互斥。RAMI4.0聚焦生产制造本身，IIRA则呈现出更强的信息技术特征；IIRA的功能视角在结构上可与RAMI4.0相对照，但RAMI4.0典型的三条轴线中有两条（生命周期与层级）在IIRA中并未被涉及。IIRA为能源、医疗卫生、制造、公共事务、交通运输等诸多领域提供了一个通用框架，而这些领域并非RAMI4.0所关照的范围；IIRA的优势还在于，其每一个视角都与其他软件参考模型保持一致，并尽可能地兼顾数据与流程的安全性。

\subsection{标准与规范现状}

工业4.0平台的组织架构下，标准化相关工作由五个工作组共同承担，分别负责参考架构、标准化与规范化，科研与创新，网络化系统的安全，法律框架条件，以及劳动与教育培训。

DIN与DKE则在参考架构、标准化与规范化这一工作组中代表德国标准化机构参与其中，其目标是形成协调一致、具备可操作性的结构，服务于所谓的数字化议程。截至目前，RAMI4.0连同工业4.0组件已经进入正式标准化进程（体现为DIN SPEC 91345），工业4.0的整体架构已基本成型，工业4.0组件的概念也已确立，而围绕各类方面与流程的子模型、工业4.0的语言以及具体实施建议等议题，则仍在持续推进之中。

具体而言，IEC通用数据字典是一套用于定义特征的标准化信息模型，它为每一项特征定义了机器可读的标识符、命名、定义以及计量单位、缩写、图形符号等其他信息；自然语言表述的概念定义主要面向人的使用，而定义中的关键信息则被抽取出来，作为属性录入这套标准化的特征信息模型之中，从而使机器处理成为可能。

电子设备描述语言（EDDL）是一种结构化的文本声明语言，在形式上类似于XML或HTML网页，它对设备参数及其相互依赖关系加以解释，为用户交互定义可视化的呈现方式，并描述系统访问设备数据的通信过程；EDDL已在IEC 61804-3与IEC 61804-4中完成标准化，其目的在于将设备开发与软件开发相互解耦，理论上可以配合任意协议使用，因而在过程自动化领域获得了广泛应用。

现场设备工具（FDT）则是一套用于统一生产设施参数化工具的构想，由FDT提供的标准化接口将工程工具与软件组件分离开来，技术人员借助被称为设备类型管理器（DTM）的软件组件对智能现场设备进行完整配置，在DTM的支持下，FDT这一构想为工程、诊断、服务与资产管理提供了集成化的功能。

自动化标记语言（AutomationML）是一种开放标准、基于XML的数据格式，用于存储和交换设施规划数据；借助AML中枢（AML.hub），数据得以在各专业领域之间关联并自动化处理，从而避免数据使用过程中的错误，并能够识别变更及潜在冲突，其中拓扑结构对应CAEX（IEC 62424）、几何与运动学对应COLLADA（Khronos Group制定）、逻辑对应PLCopen XML。eCl@ss则是一套跨行业的产品数据标准，用于对产品与服务进行分类与描述，其基础是符合规范的数据模型（DIN 4002／ISO 13584/IEC 61360），应用于采购、控制、销售、跨企业流程数据管理与工程等领域，分类工作在类别、特征与数值这几层结构中展开，并进一步细分为专业领域、主类、组、子组四个层级，其特征通过特征清单与关键词加以描述。

在通信协议层面，消息队列遥测传输协议（MQTT）是一种在数据量方面较为精简、用于信息传输的开放协议：它并不采用客户端与服务器直接通信的方式，而是由所谓的代理（Broker）与网关（Gateway）充当信息源与信息汇之间的中介，从而使得分布在不同数据网络（例如无线传感器网络）中的订阅者也能够以结构化、集中化的方式访问数据；该协议基于TCP/IP（对应OSI分层模型），与超文本传输协议（HTTP）类似，但更注重降低开销，其通信模型是一种发布—订阅模式，由发布者、网关（代理）与订阅者共同构成。

统一架构（OPC-UA）作为一项通信标准，于2006年提出、2008年正式发布，其开发目的是对OPC基金会已有十年历史的前身规范（即用于过程控制的OLE）加以整合，并在技术上予以更新；由此形成的数据模型全面、层级化且具有良好的可适配性，多个服务器可以为多个客户端提供层级化组织的信息流，通信大多经由TCP/IP实现，并可借助用户数据报协议（UDP）下的发布/订阅机制及时间敏感网络（TSN）实现实时性，在较新版本（v1.04及以后）中，OPC-UA还采用MQTT作为传输协议之一。

\subsection{方法论工具}

在参考架构与标准之外，围绕如何评估企业自身在工业4.0方面所处的水平、又该如何据此制定行动路径，业界也发展出若干成体系的方法论工具。

工业4.0就绪度模型（Readiness-Modell）是评估机械制造企业工业4.0就绪程度与能力水平的基础。该模型描绘了工业4.0的愿景，界定了企业的起点，描述了实现过程中的关键步骤与阻碍，并识别出当前领先企业所具备的特征。其基础建立在工业4.0的六个核心维度之上，即战略与组织、智能工厂、智能运营、智能产品、数据驱动服务与员工，并据此将企业划分为对应不同成熟度阶段的"新手""入门者"与"先行者"三类。

acatech工业4.0成熟度指数则支持企业规划各自专属的工业4.0发展路径，其应用过程分为三个依次递进的阶段：
\begin{enumerate}
    \item 确定成熟度，即评估企业当前的发展阶段，依据功能领域与设计维度对现有的工业4.0能力加以考察；
    \item 确定目标状态，即在此前所确定的发展阶段与目标状态基础上，借助差距分析识别出需要构建的能力；
    \item 确定行动措施，即将由此推导出的行动措施落实到路线图之中，用以构建所识别出的各项能力，图中气泡的大小即代表所达到的成熟度水平。
\end{enumerate}

VDMA 工业4.0指南则是一套面向中小企业的定位辅助工具，用于将工业4.0的整体愿景转化为企业自身专属的工业4.0构想。该指南分为准备、分析、创意、评估、导入五个阶段，以准备与分析阶段所确定的企业自身工业4.0现状为基础，进而在企业内部开展工作坊，由跨学科的企业团队以个人及小组形式共同提出并评估工业4.0构想，并将其进一步发展为具体的商业模式构想，最终将构想转化为项目，确立企业专属的工业4.0战略。工作坊所使用的工业4.0工具箱作为创意工具，分别针对产品与生产两个方向，呈现出生产与产品在不同应用层面上的发展阶段，这些发展阶段由左至右描绘出通向工业4.0愿景的路径，用以启发个性化的、贴合企业自身情况的深入思考，其创意形成过程包括确定具有潜力的应用示例、判断哪些应用层面适合进一步发展、评估该示例当前所处的发展阶段、以及探讨该示例可以有意义地推进至哪些更高发展阶段这几个步骤。


